\section{Appendix Overview}

This appendix provides implementation-independent details for the LeanNumBench
preprint. The main paper is self-contained: it states the benchmark design,
primary results, and diagnostic conclusions. The appendix records additional
dataset, protocol, verification, and analysis details so that readers can audit
the claims without relying on local project knowledge.

The appendix follows five principles. First, all claims are tied to
Lean-checked theorem records or structured experiment logs. Second, endpoint
rankings are treated as dated snapshots rather than permanent model-version
claims. Third, proof-completion passes are separated from extraction failures,
formatting failures, repair-assisted passes, and \texttt{sorry}-using outputs.
Fourth, hard-subset analysis is reported as a robustness probe, not as an
unbiased pass@3 estimate. Fifth, artifact descriptions avoid
environment-specific paths, credentials, and local machine details.

\section{Dataset Schema}

Each LeanNumBench record packages a numerical-analysis theorem and the metadata
needed to turn it into evaluation tasks. Conceptually, a record contains:

\begin{itemize}
\item an informal mathematical statement;
\item a LaTeX-style statement for human inspection;
\item a Lean theorem declaration;
\item local declarations available to the proof task;
\item provenance and source-alignment metadata when available;
\item tags for family, proof pattern, and difficulty;
\item task specifications for proof completion, statement translation, or
theorem retrieval;
\item verification metadata used by the evaluation harness.
\end{itemize}

The record is deliberately richer than a theorem string. Numerical-analysis
formalization often depends on local definitions and helper lemmas: stencil
definitions, finite-sum conventions, projection operators, local norm APIs, and
boundary bookkeeping. Preserving this context makes the task closer to the
workflow encountered by a formalizer.

\section{Family Coverage}

Table~\ref{tab:supp-family-coverage} gives the 160-record LeanNumBench v1 coverage
used by the main paper.

\begin{table}[htbp]
\centering
\caption{Family coverage in the 160-record LeanNumBench v1 release.}
\label{tab:supp-family-coverage}
\begin{tabular}{lr}
\toprule
Family & Records\\
\midrule
FEM/P0 projection & 28\\
finite-difference KdV & 25\\
spectral/DFT identities & 23\\
finite-difference heat & 15\\
linear-algebra stability & 14\\
finite-difference Burgers & 10\\
finite-difference stability & 10\\
finite-difference conservation & 8\\
interpolation & 8\\
ODE symplectic invariants & 7\\
Runge--Kutta identities & 6\\
quadrature & 6\\
\bottomrule
\end{tabular}
\end{table}

The v1 release is curated for proof-shape diversity rather than raw
count. The families stress finite sums and telescoping, CFL convexity,
boundary conditions, projection orthogonality, spectral normalization,
definition unfolding, norm inequalities, and elementary interpolation or
quadrature algebra.

\section{Provenance Policy}

The LeanNumBench v1 release contains 59 records with textbook or paper
provenance, 4 classic/general records, and 97 records seeded from the companion
Lean formalization. Provenance is used to document mathematical motivation and
to reduce the risk that the benchmark is merely a set of self-invented Lean
exercises. It does not mean that each theorem is a verbatim translation of a
textbook theorem.

For externally motivated rows, the intended provenance fields identify the
mathematical source, the numerical-analysis topic, the theorem family, and the
informal-to-formal alignment. The Lean-checked theorem remains the evaluation
object. This distinction matters because Lean formalization often requires
specialized definitions or finite-dimensional versions of textbook statements.

\subsection{Provenance-Stratified Panel}

The 160-record frozen panel contains both internally seeded records and
textbook/paper-provenance records. Aggregated over the five models, the
textbook/paper subset has 253/295 passing model-record pairs (85.8\%), while
the internal Lean-seed subset has 430/485 passing pairs (88.7\%). This
stratification is not intended as a difficulty-matched comparison, but it shows
that externally motivated records are part of the measured benchmark signal.

\begin{table}[H]
\centering
\caption{Proof-completion performance by provenance stratum in the dated
160-record frozen panel.}
\label{tab:supp-provenance-strata}
\small
\begin{tabular}{lrrrr}
\toprule
Provenance stratum & Pairs & Passed & Failed & Accuracy\\
\midrule
classic or general & 20 & 15 & 5 & 75.0\%\\
internal Lean seed & 485 & 430 & 55 & 88.7\%\\
textbook or paper & 295 & 253 & 42 & 85.8\%\\
\bottomrule
\end{tabular}
\end{table}

\section{Dataset Construction And Leakage Audit}

The May-2026 release is a development freeze rather than a hidden adaptive
leaderboard. The freeze contains 160 records, 405 task specifications, and 12
theorem families. Local verification passed 14/14 smoke suites, and the
160-row proof-prompt audit reported zero issues. The prompt audit checks that
\texttt{prompts.jsonl}, \texttt{targets.jsonl}, and
\texttt{candidate\_template.jsonl} are row-aligned; candidate templates are
empty before model generation; target proof hints are not copied verbatim into
prompts; and the source-context block before the target theorem does not
include the target theorem header.

The companion Lean library contains 176 theorem declarations. The v1 release
represents 160 of them; the remaining 16 are not a hidden test split. They are
a tracked expansion backlog, summarized in Table~\ref{tab:supp-missing-lean}.

\begin{table}[H]
\centering
\caption{Companion Lean declarations not represented in LeanNumBench v1.}
\label{tab:supp-missing-lean}
\small
\begin{tabularx}{\linewidth}{lrX}
\toprule
Priority & Count & Reason\\
\midrule
A & 10 & Interpolation, quadrature, and ODE/symplectic shape-balancing
candidates: endpoint exactness, affine linearity, or one-step map components.\\
C & 6 & Extra finite-difference KdV, Burgers, and conservation wrappers;
deferred to avoid headline expansion by near-duplicates.\\
\bottomrule
\end{tabularx}
\end{table}

Table~\ref{tab:supp-record-examples} shows three representative records. The
examples illustrate that a record includes mathematical intent, Lean-local
declarations, and a proof-completion target; it is not merely a theorem string.

\begin{table}[H]
\centering
\caption{Representative LeanNumBench record surfaces.}
\label{tab:supp-record-examples}
\scriptsize
\begin{tabularx}{\linewidth}{
  >{\raggedright\arraybackslash}p{0.24\linewidth}
  >{\raggedright\arraybackslash}p{0.24\linewidth}
  >{\raggedright\arraybackslash}p{0.24\linewidth}
  >{\raggedright\arraybackslash}X}
\toprule
Record & Informal theorem & Prompt surface / exposed formal context & Proof-completion target\\
\midrule
Lax--Friedrichs CFL stability &
Scalar Lax--Friedrichs update preserves an \(L_\infty\) bound under
\(-1 \leq \lambda \leq 1\). &
\texttt{laxFriedrichsStep}; the target theorem header is shown separately. &
Use nonnegative CFL weights, absolute-value inequalities, and arithmetic
normalization.\\
Centered finite DC mode &
Subtracting the finite mean makes the DC sum vanish. &
\texttt{finiteDC}, \texttt{finiteMean}, and \texttt{centeredSignal}; the
target theorem header is shown separately. &
Unfold definitions, rewrite finite sums of differences and constants, then
clear the nonzero-cardinality denominator.\\
P0 projection idempotence &
Averaging the two projected cell values gives the original average. &
\texttt{cellAverage2}; the target theorem header is shown separately. &
Unfold the two projection coordinates and average, then normalize by ring.\\
\bottomrule
\end{tabularx}
\end{table}

The leakage claim is intentionally narrow. Proof-completion prompts expose the
theorem statement and local declarations because that is the task. The audit
rules out copied proof hints, target-theorem duplication in the pre-target
source context, nonempty candidate templates, local path leakage, and
credential leakage; it does not claim that every local helper lemma is
semantically independent of the target proof.

\section{Neighboring Benchmarks}

Table~\ref{tab:neighboring-benchmarks} summarizes the intended boundary between
LeanNumBench and nearby benchmark families. The comparison is meant to prevent an
overbroad reading of the contribution: LeanNumBench is a diagnostic microscope for
numerical-analysis proof completion, not a claim to cover all applied
mathematics.

\begin{table}[H]
\centering
\caption{How LeanNumBench is positioned relative to neighboring benchmarks.}
\label{tab:neighboring-benchmarks}
\small
\begin{tabularx}{\linewidth}{
  >{\raggedright\arraybackslash}p{0.23\linewidth}
  >{\raggedright\arraybackslash}p{0.33\linewidth}
  >{\raggedright\arraybackslash}X}
\toprule
Benchmark family & Primary scope & LeanNumBench distinction\\
\midrule
General formal-math benchmarks & Broad theorem proving,
retrieval, and autoformalization across mathematical domains & Isolates
Lean-checked numerical-analysis proof completion and reports
domain-specific failure modes.\\
CAM-Bench-style applied-math Lean benchmarks & Broad computational and
applied mathematics at larger scale & Trades breadth for curated
numerical-analysis proof-completion families, hard-subset diagnostics,
declaration-context and repair ablations, released model-output summaries, and
proof-engineering taxonomy.\\
Construction-verification benchmarks & Tasks requiring explicit object,
algorithm, or representation construction before correctness proofs & Uses
direct proof completion over curated numerical-analysis theorem statements,
with local declaration context held fixed.\\
Formal numerical verification & Mechanized verification of numerical methods,
floating-point libraries, ODE reasoning, and PDE-related theorems & Packages
theorem-family tasks as an LLM evaluation suite rather than presenting a new
verified numerical method.\\
\bottomrule
\end{tabularx}
\end{table}

\section{Evaluation Protocol}

\subsection{Proof Completion}

The primary task is proof completion. A model receives the theorem statement
and the local declaration context. It must output only the proof body after
\texttt{:= by}. A candidate is counted as passing only if it checks in Lean
without using \texttt{sorry}.

The evaluator records several outcomes:

\begin{itemize}
\item extraction failure: no usable proof body is found in the model output;
\item rejection: the output violates task constraints before Lean checking;
\item API/transport failure: no usable model output remains after fixed retry
and materialization;
\item raw pass: the extracted proof checks without repair;
\item repair-assisted pass: a narrow deterministic repair changes a failing
candidate into a Lean-checking proof;
\item Lean failure: the candidate is well-formed enough to check but Lean
reports an error;
\item \texttt{sorry} use: the candidate uses an incomplete-proof marker and is not
counted as a pass.
\end{itemize}

\subsection{Deterministic Repair}

The repair mode is intentionally narrow. It may remove redundant no-goals tactic
lines, replace a failing \texttt{ring} step with \texttt{ring\_nf} when the
error suggests a normal-form mismatch, or drop a final tactic after Lean has
already closed all goals. It does not call a model, synthesize new proof steps,
search for lemmas, or change the theorem statement. The main paper reports raw
and repair-assisted outcomes separately.

\subsection{Statement Translation}

Statement translation asks a model to produce a Lean theorem declaration from
an informal or LaTeX-style mathematical statement. This task is useful for
calibrating whether failures are proof-specific. In the current results,
statement translation on the 50-record slice is nearly saturated, while proof
completion remains discriminating.

\subsection{Prompt And Decoding Protocol}

The one-shot proof-completion snapshot uses deterministic decoding with one
sample per model-record pair. The hard-subset pass@3 probe uses three samples
per model-record pair with nonzero sampling temperature. Empty visible proof
bodies are counted as failures, not as missing data. Transport failures and
extraction failures are tracked separately from Lean failures. Targeted retries
are permitted only to recover transport or truncation failures under the same
task definition; they do not change the theorem statement, local declaration
context, or proof-output rule.

\subsection{Model Endpoint Protocol}

The paper reports dated model snapshots rather than timeless rankings. For each
OpenRouter run we store the endpoint id, query date, temperature, output-token
cap, sample id, prompt pack, materialized candidates, and Lean-check report.
The final frozen 160-record panel uses the endpoint ids in
Table~\ref{tab:supp-model-endpoints}. Earlier 123-record snapshots are internal
pilot evidence only and are not used for the main leaderboard, hard-subset, or
failure-taxonomy claims.

\begin{table}[H]
\centering
\caption{Endpoint protocol for the final frozen proof-completion panel.}
\label{tab:supp-model-endpoints}
\small
\begin{tabularx}{\linewidth}{lXrr}
\toprule
Short name & OpenRouter endpoint id & Temp. & Max output\\
\midrule
GPT-5.5 Pro & \texttt{openai/gpt-5.5-pro} & 0.0 & 8192\\
Claude Opus 4.8 & \texttt{anthropic/claude-opus-4.8} & 0.0 & 8192\\
Gemini 3.1 Pro Preview & \texttt{google/gemini-3.1-pro-preview} & 0.0 & 8192\\
Qwen3.6 Max Preview & \texttt{qwen/qwen3.6-max-preview} & 0.0 & 8192\\
DeepSeek v4 Pro & \texttt{deepseek/deepseek-v4-pro} & 0.0 & 8192\\
\bottomrule
\end{tabularx}
\end{table}

\section{Verification Summary}

Table~\ref{tab:supp-verification-summary} summarizes the development-freeze
checks used by the main paper. These checks are reported at the level a reader
needs to audit the benchmark; repository-specific paths and local machine
details are intentionally omitted.

\begin{table}[htbp]
\centering
\caption{Development-freeze verification summary.}
\label{tab:supp-verification-summary}
\begin{tabular}{lr}
\toprule
Check & Result\\
\midrule
Benchmark records & 160\\
Task specifications & 405\\
Proof-completion tasks & 160\\
Statement-translation tasks & 160\\
Theorem-retrieval tasks & 85\\
Lean smoke suites & 14/14 passed\\
Proof-prompt rows audited & 160\\
Prompt-audit issues & 0\\
Text artifacts scanned for identity/path leaks & 193\\
\bottomrule
\end{tabular}
\end{table}

\section{Main Experiment Tables}

\subsection{Fixed Tactic Baseline}

Before calling any LLM, we run a deterministic Lean tactic portfolio over all
160 LeanNumBench v1 proof-completion rows. The portfolio contains only fixed tactic
scripts and does not use target theorem names. It solves 68/160 rows (42.5\%),
while each individual tactic solves at most 59/160 rows (36.9\%).

\begin{table}[H]
\centering
\caption{No-LLM tactic baseline over the 160-record LeanNumBench v1 release.}
\label{tab:supp-tactic-baseline}
\small
\begin{tabular}{lrrr}
\toprule
Baseline & Passed & Total & Accuracy\\
\midrule
best single fixed tactic & 59 & 160 & 36.9\%\\
fixed tactic portfolio & 68 & 160 & 42.5\%\\
\bottomrule
\end{tabular}
\end{table}

\begin{table}[H]
\centering
\caption{The 13 fixed tactic scripts used by the no-LLM portfolio.}
\label{tab:supp-tactic-scripts}
\scriptsize
\begin{tabularx}{\linewidth}{lX}
\toprule
Script ID & Proof body tried for every proof-completion row\\
\midrule
\texttt{rfl} & \texttt{rfl}\\
\texttt{simp} & \texttt{simp}\\
\texttt{simpa} & \texttt{simpa}\\
\texttt{simp\_all} & \texttt{simp\_all}\\
\texttt{norm\_num} & \texttt{norm\_num}\\
\texttt{ring} & \texttt{ring}\\
\texttt{ring\_nf} & \texttt{ring\_nf}\\
\texttt{linarith} & \texttt{linarith}\\
\texttt{nlinarith} & \texttt{nlinarith}\\
\texttt{omega} & \texttt{omega}\\
\texttt{aesop} & \texttt{aesop}\\
\texttt{simp\_ring\_nf} & \texttt{simp} followed by \texttt{ring\_nf}\\
\texttt{simp\_all\_ring\_nf} & \texttt{simp\_all} followed by \texttt{ring\_nf}\\
\bottomrule
\end{tabularx}
\end{table}

\subsection{Proof Completion Snapshot}

Table~\ref{tab:supp-proof-snapshot} reproduces the dated May-2026
five-endpoint proof-completion panel over all 160 LeanNumBench v1 records. The
table should be read as a snapshot of third-party endpoint aliases under one
prompt protocol, not as a timeless ranking.

\begin{table}[htbp]
\centering
\caption{May-2026 proof-completion panel over all 160 LeanNumBench v1 records.}
\label{tab:supp-proof-snapshot}
\begin{tabular}{lrrr}
\toprule
Endpoint alias & Passed & Total & Accuracy\\
\midrule
Gemini 3.1 Pro Preview & 147 & 160 & 91.9\%\\
GPT-5.5 Pro & 145 & 160 & 90.6\%\\
Claude Opus 4.8 & 141 & 160 & 88.1\%\\
Qwen3.6 Max Preview & 139 & 160 & 86.9\%\\
DeepSeek v4 Pro & 126 & 160 & 78.8\%\\
\bottomrule
\end{tabular}
\end{table}

\subsection{Post-Freeze Endpoint-Variant Audit}

After freezing the primary five-endpoint panel, we audited two endpoint pairs to
measure endpoint churn and cost sensitivity. The primary GPT-5.5 Pro and
Claude Opus 4.8 runs are shown beside same-prompt-pack variant runs for
GPT-5.5 standard and Claude Opus 4.7. These runs use the same 160
proof-completion prompts and local checker, but the variant rows are not used
to define the main panel, hard subset, or failure taxonomy in the main paper.

\begin{table}[htbp]
\centering
\caption{Post-freeze endpoint-variant audit on the frozen 160-record proof pack.}
\label{tab:supp-model-variant-audit}
\small
\begin{tabular}{lrrrr}
\toprule
Endpoint alias & Passed & Total & Observed run cost (USD) & Length stops\\
\midrule
GPT-5.5 & 149 & 160 & 3.55 & 0\\
GPT-5.5 Pro & 145 & 160 & 61.86 & 6\\
Claude Opus 4.7 & 134 & 160 & 2.24 & 0\\
Claude Opus 4.8 & 141 & 160 & 2.16 & 0\\
\bottomrule
\end{tabular}
\end{table}

The audit illustrates endpoint churn and cost sensitivity. Because these runs
were performed after the main panel was frozen, they are reported only as
secondary evidence that the reproducible object should be the prompt pack,
local checker, and released candidates rather than a timeless model ranking.
We therefore do not use these variants to revise the hard subset, failure
taxonomy, or headline panel. Costs are observed API-route costs for this
collection protocol and should not be interpreted as stable provider pricing.

\subsection{Statement Translation Snapshot}

On a 50-record statement-translation slice, the same model class is nearly
saturated: Gemini 3.1 Pro Preview solves 49/50 and DeepSeek v4 Pro solves
50/50 under local checking. This contrast motivates focusing the paper on
proof completion rather than only on formal statement generation. We report
this as a calibration snapshot, not as a full 160-record statement-translation
leaderboard, and we do not use it to rank models.

\subsection{Final Hard-Subset Pass@3 Probe}

After freezing the 160-record proof-completion panel, we rerun the final
11-row hard subset with three samples per model-record pair. The probe uses the
same theorem statements, declaration context, extraction rules, and Lean
checker as the one-shot panel, but samples at nonzero temperature.

\begin{table}[htbp]
\centering
\caption{Final 160-panel hard-subset pass@3 probe.}
\label{tab:supp-pass3}
\begin{tabular}{lrr}
\toprule
Model & Records solved by pass@3 & Passing samples\\
\midrule
GPT-5.5 Pro & 9/11 & 21/33\\
Gemini 3.1 Pro Preview & 9/11 & 19/33\\
DeepSeek v4 Pro & 8/11 & 13/33\\
Claude Opus 4.8 & 7/11 & 19/33\\
Qwen3.6 Max Preview & 5/11 & 10/33\\
\bottomrule
\end{tabular}
\end{table}

Pass@3 improves coverage: every hard row is solved by at least one model.
However, the hard subset is not saturated. The best model solves 9/11 rows,
and four rows remain unsolved by at least two of five models.

\begin{table}[htbp]
\centering
\caption{Row-level hard-subset persistence under pass@3.}
\label{tab:supp-hard-subset-rows}
\small
\begin{tabularx}{\linewidth}{Xrrr}
\toprule
Hard row & One-shot failures & Models solved by pass@3 & Passing samples\\
\midrule
Lax--Friedrichs CFL stability & 4/5 & 5/5 & 11/15\\
P0 best constant approximation & 4/5 & 4/5 & 8/15\\
Weighted P0 average scaling & 4/5 & 4/5 & 5/15\\
Weighted P0 residual mass zero & 4/5 & 3/5 & 5/15\\
P0 projection idempotence & 4/5 & 2/5 & 6/15\\
Centered signal zero DC mode & 3/5 & 5/5 & 10/15\\
Weighted P0 average constancy & 3/5 & 4/5 & 11/15\\
Heat CFL weight nonnegativity & 3/5 & 4/5 & 7/15\\
Finite row constant preservation & 3/5 & 4/5 & 12/15\\
KdV canonical flux const-difference & 3/5 & 2/5 & 4/15\\
Diagonal infinity-norm bound & 3/5 & 1/5 & 3/15\\
\bottomrule
\end{tabularx}
\end{table}

\section{Declaration-Context Ablation}

The declaration-context ablation tests whether local Lean declaration names are
part of the task. The full-context prompt includes local declaration names
available through imports. The no-declaration-context prompt preserves the
theorem statement, record-local definitions listed in metadata, proof-output
rules, and domain guidance, but removes the declaration-context block.

For DeepSeek v4 Pro, full declaration context solves 27/37 rows while no
declaration context solves 20/37. For Gemini 3.1 Pro Preview, full context
solves 32/37 while no-context solves 28/37. Aggregated over the two models,
full context solves 59/74 model-row pairs, while no-context solves 48/74.
Row-level flips show that context helps on 16 paired rows, no-context helps on
5, both variants pass on 43 rows, and both fail or reject on 10 rows. An
exact two-sided sign test over the 21 discordant model-row pairs gives
\(p=0.0266\), but it treats model-row pairs as independent and does not adjust
for repeated rows or repeated endpoints. We therefore use the p-value as an
exploratory descriptive summary, not as a confirmatory significance claim. The
materialized prompts average 5004.9 characters in the full-context condition
and 3217.3 characters in the no-context condition, reflecting the removed
declaration-context block.

\begin{table}[htbp]
\centering
\caption{Two-model declaration-context ablation.}
\label{tab:supp-context-ablation}
\small
\begin{tabular}{lrrrr}
\toprule
Model & Full context & No context & Hard full & Hard no-context\\
\midrule
DeepSeek v4 Pro & 27/37 & 20/37 & 4/7 & 2/7\\
Gemini 3.1 Pro Preview & 32/37 & 28/37 & 5/7 & 5/7\\
\bottomrule
\end{tabular}
\end{table}

The result should not be read as a universal law that more context always
helps. Instead, it shows that LeanNumBench measures context-sensitive proof
engineering rather than isolated theorem recall, and it motivates larger
cluster-aware context controls in future releases.

\section{Controlled Proof-Surface Pilot}

The main benchmark evaluates naturally occurring LeanNumBench records. As a
separate diagnostic, we also ran a small matched-surface pilot. Within each
pair, the two theorem targets are designed to have matched mathematical intent
while changing how much intermediate proof structure is exposed in the
statement. Each pair has an ``easy'' surface that exposes more intermediate
structure and a matched ``hard'' surface that hides the same proof obligation
behind nested APIs, index transport, or proof choreography. This pilot is not
part of the 160-record leaderboard and should be read as supplemental
mechanism evidence rather than as a new benchmark scale claim.

\begin{table}[htbp]
\centering
\caption{Controlled proof-surface pilot on 8 matched pairs.}
\label{tab:supp-controlled-pair-summary}
\small
\begin{tabular}{lrrrrrrr}
\toprule
Model & Easy & Hard & Delta & Both & Easy-only & Neither & Hard-only\\
\midrule
GPT-5.5 standard & 7/8 & 5/8 & +25.0 pp & 5 & 2 & 1 & 0\\
Claude Opus 4.8 & 6/8 & 3/8 & +37.5 pp & 3 & 3 & 2 & 0\\
\bottomrule
\end{tabular}
\end{table}

Table~\ref{tab:supp-controlled-pair-summary} shows the same aggregate
direction for two independent model families: proof surfaces with exposed
intermediate structure are easier than matched surfaces that hide the same
obligation behind more involved Lean APIs. Neither model has a hard-only pair.
The pair-level pattern is not identical across models, so the result should
not be interpreted as a target-by-target causal law.

\begin{table}[htbp]
\centering
\caption{Pair-level outcomes in the controlled proof-surface pilot.}
\label{tab:supp-controlled-pair-pairs}
\small
\begin{tabular}{llcc}
\toprule
Pair & Family & GPT-5.5 standard & Claude Opus 4.8\\
\midrule
V4P01 & heat stability & both-pass & neither-pass\\
V4P04 & Lax--Friedrichs stability & easy-only & easy-only\\
V4P06 & upwind stability & neither-pass & easy-only\\
V4P10 & KdV conservation & both-pass & easy-only\\
V4P14 & FEM/P0 projection & both-pass & both-pass\\
V4P19 & finite-row stability & easy-only & both-pass\\
V4P20 & row-stochastic stability & both-pass & neither-pass\\
V4P22 & symplectic integration & both-pass & both-pass\\
\bottomrule
\end{tabular}
\end{table}

The pilot has several caveats. It contains only 8 pairs. The GPT-5.5 standard
run has one easy-side length-truncated empty output on V4P06, while the Opus
run has no length-empty outputs but includes some fenced or explanatory model
text that must be extracted before Lean checking. The paired outcomes therefore
support a conservative claim: proof presentation and local API choreography can
change proof-completion difficulty even when the mathematical intent is
matched. They do not replace the main 160-record benchmark, the hard-subset
probe, or the failure taxonomy. The companion artifact includes the controlled
prompt pack, usage-stripped model outputs, Lean-check reports, and diagnostics under
\texttt{experiments/controlled\_pair\_v4\_selected\_8/} and
\texttt{results/controlled\_pair\_v4\_selected\_8/}.

\section{Simple-Repair Audit}

The leaderboard uses deterministic repair only where reported. On the
73-record delta slice across five models, there are 365 model-row pairs. Simple
repair changes 306 raw passes to 313 final passes, a lift of 7 pairs or 1.9
percentage points. All repair-assisted passes come from dropping redundant
no-goals tactic lines. In the final 160-record panel, only 17/800
model-record pairs are repair-assisted passes (2.1\%, Wilson 95\% CI
1.3--3.4\%). No reported pass uses \texttt{sorry}. This supports the claim
that the leaderboard is not driven by hidden post-processing.

\section{Failure Taxonomy}

The failure taxonomy classifies each non-pass outcome in the five-endpoint by
160-record frozen panel. There are 800 model-record attempts. Of these, 786
materialize as Lean proof candidates, 698 pass, 88 materialized candidates fail
Lean, and 14 attempts fail before Lean checking. The pre-Lean failures are
kept separate from Lean checker failures: 12 are empty proof bodies rejected
during model-output materialization, and 2 are API/transport failures after
the fixed collection protocol. The taxonomy below covers the 102 total
non-pass outcomes.

\begin{table}[htbp]
\centering
\caption{Failure categories for the 102 non-pass model-record attempts.}
\label{tab:supp-failure-taxonomy}
\begin{tabular}{lrr}
\toprule
Failure category & Count & Share\\
\midrule
proof-search / tactic composition & 55 & 53.9\%\\
other Lean failure & 14 & 13.7\%\\
rejected empty proof body & 12 & 11.8\%\\
type mismatch or typeclass issue & 9 & 8.8\%\\
Lean syntax / malformed proof body & 8 & 7.8\%\\
missing context or identifier & 2 & 2.0\%\\
API/transport failure & 2 & 2.0\%\\
\bottomrule
\end{tabular}
\end{table}

The dominant category is proof-search or tactic composition. This includes
unsolved goals, failed rewrites, insufficient inequality chaining, failed
arithmetic tactics, or local algebra left in the wrong normal form. Rows with
multiple errors are assigned one category by priority: malformed Lean syntax is
classified before downstream proof-search errors, and unknown identifiers are
classified before type mismatches when both occur.

\begin{table}[htbp]
\centering
\caption{Failure concentration by theorem family in the 160-record frozen panel.}
\label{tab:supp-failure-by-family}
\small
\begin{tabular}{lrrrr}
\toprule
Family & Pairs & Passed & Failed & Accuracy\\
\midrule
FEM/P0 projection & 140 & 105 & 35 & 75.0\%\\
linear-algebra stability & 70 & 53 & 17 & 75.7\%\\
spectral/DFT identities & 115 & 103 & 12 & 89.6\%\\
finite-difference KdV & 125 & 115 & 10 & 92.0\%\\
finite-difference heat & 75 & 66 & 9 & 88.0\%\\
finite-difference stability & 50 & 43 & 7 & 86.0\%\\
finite-difference conservation & 40 & 35 & 5 & 87.5\%\\
finite-difference Burgers & 50 & 46 & 4 & 92.0\%\\
ODE Runge--Kutta & 30 & 28 & 2 & 93.3\%\\
ODE symplectic invariants & 35 & 34 & 1 & 97.1\%\\
interpolation & 40 & 40 & 0 & 100.0\%\\
quadrature rules & 30 & 30 & 0 & 100.0\%\\
\bottomrule
\end{tabular}
\end{table}

\subsection{Classification Rules}

Each non-pass model-record attempt receives one primary category. The
categories are assigned by priority so that a single attempt contributes one
count. Pre-check rejections are separated before Lean-error categories.
Malformed Lean syntax is classified before downstream proof-search errors;
missing identifiers are classified before type mismatches; otherwise the
remaining unsolved proof obligations are classified as proof-search or
tactic-composition failures. The fallback category is reserved for rare Lean
failures that do not fit the preceding classes. The current release uses one
fixed-priority adjudication pass rather than independent double annotation;
row-level result summaries and usage-stripped model outputs are included for
audit.

\begin{table}[htbp]
\centering
\caption{Failure-taxonomy categories used in the main paper.}
\label{tab:supp-taxonomy-rules}
\small
\begin{tabularx}{\linewidth}{lX}
\toprule
Category & Typical evidence\\
\midrule
Rejected empty proof body & The model output contains no nonempty proof body
after task-format extraction and is rejected before Lean checking.\\
API/transport failure & No usable model output remains after fixed retry and
materialization.\\
Lean syntax / malformed proof body & Unexpected tokens, malformed tactic
blocks, or proof text that Lean cannot parse.\\
Missing context or identifier & Unknown declaration, namespace, theorem name,
or local identifier.\\
Type mismatch or typeclass issue & Application mismatch, incompatible expected
type, or failed typeclass synthesis.\\
Proof-search / tactic composition & Parsed proof leaves unsolved goals, failed
rewrites, arithmetic side conditions, no-progress tactics, or normal-form
mismatches.\\
Other Lean failure & Rare residual failures outside the previous categories.\\
\bottomrule
\end{tabularx}
\end{table}

\paragraph{Released Lean-error excerpts.}
The following ASCII-normalized excerpts are short checker messages from the
released controlled-pair diagnostics. They illustrate the evidence used by
the taxonomy without exposing local filesystem paths.

\begin{verbatim}
Proof-search:
Candidate.lean:32:2: error: unsolved goals
goal: -1 <= cfl
\end{verbatim}

\begin{verbatim}
Application/type-shape mismatch:
error: Tactic `apply` failed: could not unify
  |heatStepGrid r u i| <= M
with
  |heatStepGrid r (fun j => heatStepGrid r u j) i| <= M
\end{verbatim}

\begin{verbatim}
Syntax / malformed proof body:
Candidate.lean:32:3: error: unexpected token 'fun';
expected '{' or tactic
\end{verbatim}

\begin{verbatim}
Type mismatch / typeclass:
error: Application type mismatch
error(lean.synthInstanceFailed): failed to synthesize
  AddGroup Prop
\end{verbatim}

\section{Hard-Subset Case Studies}

\subsection{Lax--Friedrichs CFL L-Infinity Stability}

The row \texttt{finite\_difference.lax\_friedrichs\_cfl\_linf\_stability\_1d}
is failed by four of five models in one-shot evaluation. Pass@3 raises it to
five of five models, but only 11 of 15 samples pass. The proof is a CFL
convex-combination argument: use nonnegativity of the Lax--Friedrichs stencil
weights, distribute absolute values across weighted terms, combine max bounds,
and close arithmetic. In Lean, this becomes a chain of local lemma selection,
absolute-value rewriting, and inequality side conditions.

\subsection{Centered Signal Has Zero DC Mode}

The general-\(N\) centered-DC row states that subtracting the finite mean makes
the DC sum vanish. The Lean proof requires unfolding finite mean and DC
definitions, rewriting a finite sum of differences, rewriting a constant sum
over a finite set, using the nonzero-cardinality assumption, and closing the
normalized algebra. Several model outputs reach a nearly algebraic remaining
goal but fail to align the finite-sum rewrites.

\subsection{P0 Projection Idempotence}

The two-cell P0 projection-idempotence row is mathematically elementary:
averaging projected cell values gives the same average. Yet multiple models
unfold only part of the local projection API or leave Lean with a normal-form
mismatch. This row illustrates why short rows can still be diagnostic:
numerical-analysis formalization often fails on local definition unfolding and
normalization details that informal mathematics hides.

\subsection{Diagonal Infinity-Norm Bound}

A backup hard case is the diagonal infinity-norm bound. The proof requires
unfolding a small local norm API, splitting a \texttt{max\_le} goal, and
applying monotonicity of multiplication with correct side conditions. This row
also illustrates a non-monotone context effect in the ablation: the no-context
variant can pass while the full-context variant fails, showing that more
declarations can sometimes distract a model into overcomplicated proof paths.

\section{Artifact Release Package}

The reproducibility artifact contains enough material to audit the
reported tables and rerun the local Lean checks without relying on local
author paths. The LeanNumBench v1 release package includes a reproducibility
artifact zip with SHA256
\texttt{a576d3c7b7df0e37}\allowbreak\texttt{1705dbf9e4a9ec78}\allowbreak\texttt{8f9719c4b4f8ff48}\allowbreak\texttt{8483db1a9452d130}.
The package contains the components in the checklist below.

\begin{center}
\small
\textbf{Reproducibility artifact checklist.}

\vspace{0.3em}
\begin{tabularx}{\linewidth}{lX}
\toprule
Artifact component & Purpose\\
\midrule
Benchmark metadata & Defines records, families, tags, provenance, and task
types.\\
Companion Lean source & Provides theorem declarations and local
definitions used by proof candidates.\\
Prompt templates and frozen prompt packs & Fixes the input distribution for
reported model runs.\\
Model-output summaries & Supports reconstruction of reported pass, failure,
repair, and extraction counts.\\
Lean error diagnostics & Provides released checker excerpts for controlled
proof-surface failures and row-level failure summaries for the frozen
160-record panel.\\
Reproduction scripts & Provide offline artifact smoke checks, prompt audits,
model-output materialization, Lean candidate checking, deterministic repair
boundaries, and no-LLM tactic baselines.\\
Generated tables and figure data & Allows readers to trace paper figures and
tables back to structured experiment summaries.\\
Identity/path audit & Checks that released text artifacts do not expose author
identity, credentials, or environment-specific paths.\\
\bottomrule
\end{tabularx}
\end{center}

The artifact README includes exact no-API commands. The fastest artifact check
is a metadata, hash, and identity/path smoke test:

\begin{verbatim}
python reproduction_scripts/artifact_smoke_check.py
\end{verbatim}

To rerun Lean checking, the reader can build the companion Lean project and
materialize candidates from the included usage-stripped model outputs. The
artifact README gives the concrete package target for the build command:

\begin{verbatim}
cd companion_lean
lake exe cache get
lake build LeanNumerics
cd ..
RS=reproduction_scripts
OUT=_repro/leannumbench160_candidates
python $RS/materialize_model_pilot_candidates.py \
  results/raw_outputs/frontier_160_no_usage.jsonl \
  --targets experiments/model_pilot_160_proof/targets.jsonl \
  --output-dir $OUT
python $RS/check_lean_candidates_divide.py \
  --candidate-file $OUT/proof_candidates.jsonl \
  --lean-project companion_lean --repair-simple \
  --output _repro/leannumbench160_check.json
\end{verbatim}

These commands use only released text files, the included Lean source, and the
public Lake/Mathlib dependency lockfile. The cache-get step only accelerates
dependency setup; if unavailable, Lake can build from source. The commands do
not call model APIs. The expected replay materializes 786 proof candidates and
reproduces the frozen per-model pass counts in
Table~\ref{tab:supp-proof-snapshot}. The companion Lean project is pinned to
\texttt{leanprover/lean4:v4.29.1} with Mathlib revision \texttt{v4.29.1}. The
included \texttt{lake-manifest.json} SHA256 is:
\begin{verbatim}
b5199815f98f54813f7ff29fffebd34f0f421d16fd0eeacacd198f20cd7d73c1
\end{verbatim}
The released companion Lean source has also been scanned for
benchmark-relevant \texttt{sorry}, \texttt{admit}, and user-defined
\texttt{axiom} declarations, with zero matches.

The benchmark metadata are released under CC-BY-4.0. The companion Lean
source, evaluation scripts, and release tooling are released under
Apache-2.0.

\section{Versioned Benchmark Scope}

LeanNumBench v1 freezes 160 records; future versions may expand the record
count, but the paper treats v1 as the diagnostic release evaluated here rather
than a claim to out-scale broad applied-mathematics benchmarks. The
statement-translation result is a 50-record calibration slice, not a full
secondary leaderboard. The declaration-context ablation covers two models, so
it is diagnostic rather than exhaustive. The hard subset is selected post-hoc
from the one-shot leaderboard. Many records are internally seeded from the
companion Lean library, although textbook and paper provenance improves
auditability. The failure taxonomy uses fixed priority rules and is released
for row-level audit, but it is not yet backed by an independent
inter-annotator reliability study.

These limitations are intentionally exposed because LeanNumBench is meant to
support auditable rolling evaluation. The next versioned benchmark steps are a
larger pre-registered hard split, Lean-feedback and retrieval-augmented prover
baselines, a full-160 pass@k panel, shuffled or irrelevant-context controls,
and a direct comparison with numerical-analysis subsets of broader
applied-mathematics Lean benchmarks. The v1 release freezes the model panel,
reruns the leaderboard on all 160 records, regenerates the failure taxonomy
for the frozen panel, and packages a reproducibility artifact with
path-sanitized files and usage-stripped model outputs.
